\chapter{Evaluation} \label{sec:eval}
In this chapter, the test methodology, the different test scenarios and their results are presented.  
Chapter \ref{subsec:procedure} discusses the boundary conditions of the test runs.  
Subsequently, in Chapters \ref{subsec:low}, \ref{subsec:medium} and \ref{subsec:high}, the procedure is applied under environments with varying levels of interference.  
Finally, Chapter \ref{subsec:disc} compares the results of this thesis with findings from related work and places them in context.  

\section{Test procedure} \label{subsec:procedure}

Certain factors remain constant across all test runs.  
The test procedure is defined by five key elements:  
\begin{itemize}
    \item The \textbf{cycle time}, which in this thesis is fixed at 3 ms as previously described.  
    \item The \textbf{packet size of bulk data transmission}, which must be determined to minimize interference with the real-time application.  
    \item The \textbf{duration of a single test run}, i.e., the time during which the RT-application executes its computation.  
    \item The \textbf{number of repetitions} of each test run, i.e., how often a given test case is executed to produce sufficient data for reliable conclusions.  
    \item The \textbf{different test cases}, derived from the Wi-Fi factors outlined in Chapter \ref{sec:implementation}.  
    In this thesis, a total of 112 different test cases are considered.  
\end{itemize}  

To determine the execution time of the real-time application, both the number of repetitions and potential clock drift must be taken into account.  
In this thesis, each test case is executed for 60 seconds.  
As described in \cite[p. 3]{drift}, quartz clocks can exhibit a drift of 1-100 ppm (parts per million).  
This corresponds to a drift of 60 $\mu$s to 6 ms per minute.  
However, in the results of this thesis, no noticeable drift was observed within a one-minute interval.  

With a cycle time of 3 ms, each 60-second run consists of 20,000 cycles.  
Each test case is repeated three times, resulting in 60,000 cycles per case to increase the probability of capturing anomalies.  
Additionally, before each run, the system clock is synchronized, effectively eliminating clock drift.  
The Wi-Fi connection is also given a warm-up period of 60 seconds, as recommended in \cite[Chap.~3.3.3.14]{alexander-bmwg-wlan-switch-meth-01}, to avoid distortions in the results.  

The optimal packet size, as shown in Table~X, is 4096 bytes per packet.  
At this size, both latency and jitter are significantly lower than with 512-byte packets or the standard size of 1440 bytes.  

The test cases are defined as combinations of Wi-Fi standard, frequency, bandwidth, \ac{QoS}, and transport protocol.  
For Wi-Fi~4, three combinations exist: the 2.4 GHz band and in the 5 GHz band both 20 MHz and 40 MHz channel widths.  
For Wi-Fi~5, the 5 GHz band with 20 MHz, 40 MHz and 80 MHz channel widths yields three combinations.  
Wi-Fi~6 supports all three frequency bands.  
In the 2.4 GHz and 5 GHz bands, four combinations arise, while in the 6 GHz band, the available channel widths of 20, 40, 80 and 160 MHz yield four more combinations.  
In total, Wi-Fi~6 therefore contributes eight combinations.  
Together, across all standards, 14 unique combinations of standards, frequency bands, and bandwidths are considered.  

Each of these 14 combinations is tested with four differnet levels of \ac{QoS} and with both \ac{TCP} and \ac{UDP}.  
This results in a total of $8 \times 14 = 112$ test cases.  


\section{Low disturbance test} \label{subsec:low}
\section{Average disturbance test}  \label{subsec:medium}
\section{High disturbance test} \label{subsec:high}
\section{Discussion} \label{subsec:disc}